\section{Introduction}

Global warming continues to worsen each year, driven by rising levels of CO$_2$ and other greenhouse gases in the atmosphere \cite{ipcc2023}. One of the most effective ways to combat climate change is by reducing emissions from the energy sector, which is the largest contributor to global carbon output \cite{iea2022}. A central strategy for this transition is the expansion of renewable energy sources, with photovoltaics (PV) emerging as one of the most widely adopted technologies over the past decades \cite{fraunhofer2021}.

However, PV systems face a key challenge: their efficiency depends heavily on solar radiation. Since solar irradiance is strongly influenced by weather conditions, cloud cover in particular causes significant fluctuations in power output. These fluctuations can compromise grid stability, making accurate forecasting of PV generation essential \cite{lorenz2009}. Tracking cloud movement, therefore, plays a critical role in predicting PV output.

Traditionally, cloud motion is derived from satellite imagery, which is widely used for large-scale weather forecasting on the scale of hours to days \cite{schmetz2002}. While valuable for meteorology, this approach is less suitable for PV applications due to two main limitations: low spatial resolution, which prevents the detection of smaller cloud structures, and low temporal resolution, as satellites typically provide images at relatively infrequent intervals \cite{hammer2003,inman2013}.

An established alternative for long-term forecasting is \textit{nowcasting} using all-sky imagers (ASIs) \cite{chow2011,kleissl2013}. These instruments capture a hemispheric view of the sky at intervals of less than a minute, providing high-resolution observations of cloud formation, deformation, and movement. While ASIs are well-suited for this purpose, extracting accurate cloud motion vectors (CMVs) from the imagery remains a challenge. Over time, various methods have been developed, with \textbf{optical flow (OF)} techniques proving to be among the most reliable \cite{urquhart2015}.

Optical flow algorithms estimate the apparent motion of pixels between consecutive frames by analyzing changes in pixel intensity \cite{horn1981determining}. This approach eliminates the need for explicit cloud segmentation or feature tracking, making it highly effective for handling the complex and irregular motion of clouds.

Optical flow can be implemented in different ways:

\begin{itemize}

\item \textbf{Dense optical flow} (e.g., Farnebäck, Horn--Schunck) calculates motion vectors for every pixel, producing detailed motion fields but at a high computational cost \cite{farneback2003two,horn1981determining}.

\item \textbf{Sparse optical flow} (e.g., Lucas--Kanade) focuses only on selected pixels with significant features, requiring less computation and making it widely used in applications such as object tracking \cite{lucas1981iterative}.

\end{itemize}

More recently, deep learning--based approaches have been introduced to enhance CMV estimation \cite{doshi2022}. These models address challenges such as identifying cloud types, estimating cloud height, and handling multi-layered cloud systems. They also demonstrate improved robustness to common sources of error in ASI data, including image noise, brightness variations, and fisheye lens distortion \cite{shi2015,nie2024}.

In summary, the combination of all-sky imagers and optical flow techniques currently represents the most effective strategy for PV nowcasting. Ongoing research aims to further improve the robustness of CMV estimation by mitigating noise and brightness fluctuations, as well as enhancing the quality of ASI imagery \cite{jung2020}.
